% !TeX program = lualatex
% Compilar desde esta carpeta: latexmk -lualatex -interaction=nonstopmode -halt-on-error Plan_trabajo_12_semanas.tex
\documentclass[11pt,a4paper,oneside]{report}
\usepackage{fontspec}
\setmainfont{Arial}
\usepackage[spanish,es-nodecimaldot,es-noshorthands]{babel}
\usepackage{geometry}
\geometry{a4paper,left=3cm,top=3cm,right=2.5cm,bottom=2.5cm,headheight=14pt,headsep=0.6cm}
\usepackage{setspace}
\onehalfspacing
\usepackage{graphicx}
\usepackage{booktabs}
\usepackage{longtable}
\usepackage{tabularx}
\usepackage{array}
\usepackage{pdflscape}
\usepackage{ragged2e}
\usepackage{xcolor}
\usepackage{colortbl}
\usepackage{fancyhdr}
\usepackage{titlesec}
\usepackage[hidelinks]{hyperref}
\usepackage{bookmark}
\usepackage{enumitem}
\usepackage{microtype}
\usepackage{tikz}
\usetikzlibrary{arrows.meta,positioning}
\usepackage{pifont}
\definecolor{unired}{HTML}{8B1A1A}
\definecolor{darkblue}{HTML}{17365D}
\definecolor{lightblue}{HTML}{DCE6F1}
\definecolor{baseline}{HTML}{5A7D71}
\definecolor{active}{HTML}{E9B44C}
\definecolor{planned}{HTML}{DCE6F1}
\setlength{\parindent}{0pt}
\setlength{\parskip}{9pt}
\setlength{\emergencystretch}{3em}
\setlist{leftmargin=1.1cm,itemsep=0.35\baselineskip,topsep=0.35\baselineskip}
\titleformat{\chapter}[hang]{\normalfont\bfseries\fontsize{11}{13.2}\selectfont}{\MakeUppercase{\chaptername}\ \thechapter.}{0.75em}{\MakeUppercase}
\titleformat{name=\chapter,numberless}[block]{\normalfont\bfseries\fontsize{11}{13.2}\selectfont}{}{0pt}{\MakeUppercase}
\titleformat{\section}{\normalfont\bfseries\fontsize{11}{13.2}\selectfont}{\thesection}{0.75em}{\MakeUppercase}
\titleformat{\subsection}{\normalfont\bfseries\fontsize{11}{13.2}\selectfont}{\thesubsection}{0.75em}{\MakeUppercase}
\titlespacing*{\chapter}{0pt}{0pt}{0.75\baselineskip}
\fancypagestyle{plain}{\fancyhf{}\fancyhead[R]{\thepage}\renewcommand{\headrulewidth}{0pt}}
\pagestyle{plain}
\newcommand{\autores}{LUIS ENRIQUE KOC G\'ONGORA\\HERBERT ANTONIO MEL\'ENDEZ GARC\'IA}
\newcommand{\asesor}{Samuel Oporto}
\newcommand{\caratula}[1]{%
  \hypersetup{pageanchor=false}
  \begin{titlepage}
  \thispagestyle{empty}
  \begin{center}
  {\bfseries\fontsize{19}{23}\selectfont UNIVERSIDAD NACIONAL DE INGENIER\'IA\par}
  \vspace{0.2cm}
  {\bfseries\fontsize{13}{16}\selectfont UNIDAD DE POSGRADO DE LA FACULTAD DE INGENIER\'IA INDUSTRIAL Y DE SISTEMAS\par}
  {\bfseries\fontsize{13}{16}\selectfont ESCUELA DE POSGRADO\par}
  \vspace{0.25cm}
  \includegraphics[width=4.6cm]{../Plan/imagenes/logo_UNI.png}\par
  \vspace{0.2cm}
  {\bfseries\fontsize{14}{17}\selectfont #1\par}
  \vspace{0.35cm}
  {\bfseries\fontsize{12}{15}\selectfont DISE\~NO Y EVALUACI\'ON DE UNA RED NEURONAL INFORMADA POR F\'ISICA PARA ESTIMAR LA TEMPERATURA Y LA AMPACIDAD DE CABLES EL\'ECTRICOS ENTERRADOS EN ENTORNOS HETEROG\'ENEOS\par}
  \vfill
  {\fontsize{11}{14}\selectfont PARA OBTENER EL GRADO ACAD\'EMICO DE\\MAESTRO EN CIENCIAS CON MENCI\'ON EN INTELIGENCIA ARTIFICIAL\par}
  \vspace{0.35cm}
  {\bfseries\fontsize{11}{14}\selectfont PRESENTADO POR:\par}
  {\bfseries\fontsize{11}{14}\selectfont \autores\par}
  \vspace{0.3cm}
  {\fontsize{11}{14}\selectfont ASESOR: \asesor\par}
  \vfill
  {\fontsize{11}{14}\selectfont Lima, Per\'u\\2026\par}
  \end{center}
  \end{titlepage}
  \hypersetup{pageanchor=true}
  \pagenumbering{arabic}
  \tableofcontents
  \clearpage
}
\begin{document}
\caratula{PLAN DE TRABAJO PARA REVISAR Y COMPLETAR EL PLAN DE TESIS}

\chapter*{Prop\'osito y alcance}
\addcontentsline{toc}{chapter}{Prop\'osito y alcance}
Este documento organiza doce semanas de trabajo para revisar y completar el plan de tesis profesionalizante de acuerdo con la Gu\'ia Metodol\'ogica N.\textsuperscript{o} 01 de la Universidad Nacional de Ingenier\'ia. Su cronograma administra exclusivamente las tareas de revisi\'on y cierre del plan de tesis. No sustituye ni modifica el cronograma de ejecuci\'on de la tesis, que en el documento vigente abarca julio de 2026 a junio de 2027.

\begin{tabularx}{\textwidth}{>{\bfseries}p{4.4cm}X}
Inicio del periodo & 15 de septiembre de 2026 \\
Corte del avance & 10 de octubre de 2026 \\
Cierre de doce semanas & 7 de diciembre de 2026 (ultimo d\'ia incluido) \\
Limite de 84 d\'ias & 8 de diciembre de 2026 \\
Responsables & Luis Enrique Koc G\'ongora y Herbert Antonio Mel\'endez Garc\'ia \\
Asesor & Samuel Oporto \\
\end{tabularx}

\section*{Criterios de cumplimiento de la Gu\'ia 1}
La revisi\'on debe comprobar los datos generales, la car\'atula, el prop\'osito de graduaci\'on, el nombre completo de ambos autores y del asesor. Para la modalidad profesionalizante, la estructura de referencia contiene cinco cap\'itulos: (I) Planteamiento del problema, con diagn\'ostico, antecedentes bibliogr\'aficos, formulaci\'on, justificaci\'on y alcances; (II) Objetivos; (III) Marco te\'orico y definici\'on de t\'erminos; (IV) Metodolog\'ia, con tipo, nivel, m\'etodos, poblaci\'on y muestra cuando aplique, dise\~no, recolecci\'on, an\'alisis/procesamiento y etapas; y (V) Administraci\'on, con cronograma, presupuesto y financiamiento. Se revisar\'an adem\'as referencias en estilo APA vigente y anexos pertinentes. La matriz de consistencia no se asumir\'a obligatoria para esta modalidad por la Gu\'ia 1; cualquier anexo adicional se justificar\'a por necesidad del plan.

\section*{Avance documentado al 10 de octubre}
Al corte se encuentran disponibles el plan vigente y su PDF rev. 2, un informe de auditor\'ia fechado el 29 de septiembre de 2026, un inventario previo de 16 instancias de datos (nueve f\'isicas o de campo y siete de verificaci\'on), y un cat\'alogo PINN\_v2 con 21 casos activos que el cargador actual acepta. El informe previo ya incluye CIGRE TB 880, TB 963 y el caso anal\'itico annulus. Estas existencias constituyen una l\'inea base comprobable; los archivos no permiten atribuir con certeza cada producto a una semana espec\'ifica del periodo. Por ello, las semanas 1--3 se muestran como avance documental disponible, no como una declaraci\'on de fechas de ejecuci\'on verificadas. La semana 4 est\'a en curso.

\section*{Entregables de cierre}
\begin{enumerate}
  \item Matriz de cumplimiento de Gu\'ia 1 y lista priorizada de observaciones del plan profesionalizante.
  \item Versiones corregidas del texto del plan y sus fuentes, con revisi\'on de coherencia entre problema, objetivos, marco, m\'etodo y administraci\'on.
  \item Informe de datos y libro Excel navegable, con fuentes, criterios de inclusi\'on, procedencia, procesamiento y datos faltantes.
  \item Paquete JSON documentado con la versi\'on vigente de PINN\_v2; casos ejecutables y candidatos incompletos separados.
  \item Propuesta de al menos dos revistas con art\'iculos comparables y verificaci\'on editorial actual.
  \item Entrega reproducible: fuentes LaTeX, PDF compilado, libros Excel, ZIP JSON y manifiesto de validaciones.
\end{enumerate}

\begin{landscape}
\section*{Cronograma de trabajo del plan de tesis}
\addcontentsline{toc}{section}{Cronograma de trabajo del plan de tesis}
\small
\setlength{\tabcolsep}{3pt}
\renewcommand{\arraystretch}{1.35}
\begin{tabularx}{\linewidth}{p{0.65cm}X*{12}{>{\centering\arraybackslash}p{1.05cm}}}
\toprule
\textbf{ID} & \textbf{Actividad / producto} & \textbf{S1} & \textbf{S2} & \textbf{S3} & \textbf{S4} & \textbf{S5} & \textbf{S6} & \textbf{S7} & \textbf{S8} & \textbf{S9} & \textbf{S10} & \textbf{S11} & \textbf{S12} \\
\midrule
1 & Inventariar plan vigente, Gu\'ia 1 y observaciones previas & \cellcolor{baseline} & \cellcolor{baseline} & \cellcolor{baseline} & \cellcolor{active} & & & & & & & & \\
2 & Verificar estructura profesionalizante, datos generales y APA & \cellcolor{baseline} & \cellcolor{baseline} & \cellcolor{baseline} & \cellcolor{active} & \cellcolor{planned} & & & & & & & \\
3 & Revisar coherencia de diagn\'ostico, preguntas, objetivos y metodolog\'ia & & \cellcolor{baseline} & \cellcolor{baseline} & \cellcolor{active} & \cellcolor{planned} & \cellcolor{planned} & & & & & & \\
4 & Consolidar inventario de datos, campos y procedencia & & \cellcolor{baseline} & \cellcolor{baseline} & \cellcolor{active} & \cellcolor{planned} & \cellcolor{planned} & \cellcolor{planned} & & & & & \\
5 & Buscar proyectos COES y activos MINEM de cable subterr\'aneo, cualquier tensi\'on & & & & \cellcolor{active} & \cellcolor{planned} & \cellcolor{planned} & \cellcolor{planned} & & & & & \\
6 & Contrastar Arias, otros casos peruanos y CIGRE TB 880/TB 963 & & & & & \cellcolor{planned} & \cellcolor{planned} & \cellcolor{planned} & \cellcolor{planned} & & & & \\
7 & Reconciliar datos y estados de procedencia con PINN\_v2 & & & & & & \cellcolor{planned} & \cellcolor{planned} & \cellcolor{planned} & \cellcolor{planned} & & & \\
8 & Redactar informe de datos y construir libro Excel navegable & & & & & & & \cellcolor{planned} & \cellcolor{planned} & \cellcolor{planned} & \cellcolor{planned} & & \\
9 & Preparar paquete JSON, candidatos faltantes y manifiesto & & & & & & \cellcolor{planned} & \cellcolor{planned} & \cellcolor{planned} & \cellcolor{planned} & & & \\
10 & Evaluar revistas y redactar propuesta justificada & & & & & & & & & \cellcolor{planned} & \cellcolor{planned} & \cellcolor{planned} & \\
11 & Incorporar comentarios del asesor y corregir el plan & & & & & & & & & & \cellcolor{planned} & \cellcolor{planned} & \cellcolor{planned} \\
12 & Compilar, validar consistencia y cerrar los entregables & & & & & & & & & & & & \cellcolor{planned} \\
\bottomrule
\end{tabularx}
\vspace{0.35cm}

\scriptsize
\begin{tabularx}{\linewidth}{p{2.2cm}X}
\cellcolor{baseline}\rule{0.4cm}{0.25cm} & Productos o fuentes que ya existen al corte; no certifica la fecha semanal de ejecuci\'on. \\
\cellcolor{active}\rule{0.4cm}{0.25cm} & Semana en curso al 10 de octubre de 2026. \\
\cellcolor{planned}\rule{0.4cm}{0.25cm} & Trabajo programado; fechas sujetas a coordinaci\'on con el asesor. \\
\end{tabularx}

\vspace{0.25cm}
\textbf{Fechas por semana:} S1 15--21 sep.; S2 22--28 sep.; S3 29 sep.--5 oct.; S4 6--12 oct.; S5 13--19 oct.; S6 20--26 oct.; S7 27 oct.--2 nov.; S8 3--9 nov.; S9 10--16 nov.; S10 17--23 nov.; S11 24--30 nov.; S12 1--7 dic. de 2026.
\end{landscape}

\chapter*{Reglas de manejo y aceptaci\'on}
\addcontentsline{toc}{chapter}{Reglas de manejo y aceptaci\'on}
\begin{itemize}
  \item El cronograma de este documento es un plan de trabajo para completar y revisar el plan de tesis. El cronograma de investigaci\'on aprobado permanece en el Cap\'itulo V del plan vigente.
  \item Cada dato debe conservar fuente, ubicaci\'on (p\'agina/tabla), unidad, procedimiento de extracci\'on y estado de procedencia. Distinguir casos de instalaci\'on real, proyectos planificados, benchmarks anal\'iticos/FEM y escenarios de verificaci\'on.
  \item La b\'usqueda COES incluye anteproyectos y planes de transmisi\'on de cualquier tensi\'on. Distinguir propuesta de activo operativo; la tensi\'on, longitud o ampacidad no completan por s\'i solas la geometr\'ia de zanja ni las propiedades del suelo.
  \item El esquema activo PINN\_v2 es \texttt{pinn-v2-case-1}. El validador rechaza grupos requeridos marcados como faltantes; esos registros ir\'an en el ZIP como candidatos no ejecutables, sin valores geom\'etricos inventados ni cambios al esquema.
  \item Los productos finales se consideran cerrados cuando las fuentes coinciden con informe, Excel y JSON; compilan los documentos LaTeX; los casos ejecutables pasan el cargador PINN\_v2; y las limitaciones est\'an expuestas.
\end{itemize}

\chapter*{Fuentes de control}
\addcontentsline{toc}{chapter}{Fuentes de control}
Universidad Nacional de Ingenier\'ia, Escuela de Posgrado. (2023). \textit{Gu\'ia metodol\'ogica para la elaboraci\'on y presentaci\'on del Plan de Tesis (Gu\'ia N.\textsuperscript{o} 01)}. Secci\'on 5.1, Anexo III y Anexo VIII.\\
Plan de tesis profesionalizante vigente del Grupo 91, fuente LaTeX y PDF rev. 2.\\
Auditor\'ia general del plan de tesis, 29 de septiembre de 2026.
\end{document}